\documentclass[10pt,a4paper]{article}
\usepackage[margin=1in]{geometry}
\usepackage[utf8]{inputenc}
\usepackage[T1]{fontenc}
\usepackage{amsmath, amssymb, amsfonts}
\usepackage{enumitem}
\usepackage{booktabs}
\usepackage{cite}

% Define a custom command for notes/comments
\newcommand{\note}[1]{\textit{\textcolor{darkgray}{#1}}}

\begin{document}

\section{Experimental Design and Matrix}

This section outlines the \textbf{Experimental Design and Matrix} exclusively. The design is strictly centered around the primary pipeline: \textbf{Gen1 Event Detection + EAS-SNN Baseline + HAPQ (NAS + Structured Pruning + Weight/Membrane Quantization) + FPGA Deployment Verification}. The details for EAS-SNN and its Gen1/1Mpx settings are directly aligned with its original implementation and protocols.

\subsection{A. Gen1: Comparative and Ablation Experiments (Goal: Accuracy Stability)}

\subsubsection{A1) Group Settings (Accuracy Bench)}

\textbf{Fixed Settings:} Identical data split/evaluation scripts; identical event representation and sampling strategy; identical training epochs and augmentation (aligned with EAS-SNN public settings [1]).

\paragraph{SNN Mainline (Core)}
\begin{enumerate}
    \item \textbf{SNN-FP}: EAS-SNN Original (Full Precision) [1].
    \item \textbf{SNN-P}: Structured Pruning only (Block/Channel levels) [48, 78].
    \item \textbf{SNN-Q(W)}: Weight Quantization only (QAT; Mixed/Uniform) [43, 42, 38, 39].
    \item \textbf{SNN-Q(W+U)}: Weight + Membrane Potential Quantization (QAT) [76, 77].
    \item \textbf{SNN-HAPQ}: NAS-selected Structure + Structured Pruning + W/U Quantization (Proposed Method) [38--41, 78, 76, 77].
\end{enumerate}

\paragraph{ANN Control Group (Representative Models, Mandatory)}
\begin{itemize}
    \item \textbf{ANN-Transformer/RNN}: RVT-like event detection (Strong Baseline) [8].
    \item \textbf{ANN-Memory/RNN}: ASTMNet-like continuous event detection (Strong Baseline) [11].
    \item \textit{Optional Supplement}: GET-like Transformer detection [10] (Control for ``Non-RNN Attention'').
\end{itemize}

\begin{quote}
    \textit{Note: The ANN group serves to demonstrate ``SNN accuracy stability after compression'' while providing an ``upper bound/reference for mature ANN approaches.''}
\end{quote}

\subsubsection{A2) Metrics and Statistical Protocol (Gen1)}
\begin{itemize}
    \item \textbf{Primary Metrics}: mAP$_{50:95}$, mAP$_{50}$ (Consistent with EAS-SNN Gen1 tables) [1].
    \item \textbf{Stability Metrics}: mAP Variance ($\ge 3$ seeds); Convergence curves (Training/Validation).
    \item \textbf{Sparsity/Complexity Metrics}: Params, SynOps/AC Ops (SNN activity-related) [1].
    \item \textbf{Latency Proxy (Software Only)}: Inference time per sample (GPU/CPU), used only for trend reference (not as FPGA conclusions).
\end{itemize}

\subsubsection{A3) Ablation Grid}
\textit{Note: Only one variable is modified per dimension, while others remain fixed in the SNN-FP $\to$ SNN-HAPQ training/fine-tuning pipeline.}

\begin{enumerate}
    \item \textbf{Pruning Granularity}
    \begin{itemize}
        \item Channel pruning vs. SIMD-block pruning (Hardware Alignment) [78, 48].
        \item Block size (Aligned with DSP/BRAM packing width vs. Misaligned).
    \end{itemize}
    
    \item \textbf{Pruning Strength}
    \begin{itemize}
        \item Retention rate $\kappa \in \{1.0, 0.8, 0.6, 0.4\}$ (At least 4 levels).
        \item Comparison: Different structures under equivalent Params (Proving the necessity of ``structural alignment'') [48].
    \end{itemize}
    
    \item \textbf{Quantization Objects}
    \begin{itemize}
        \item W only vs. W+U (Membrane Potential/State variables) [76, 77].
        \item Bitwidth of $U$: $b_U \in \{16, 12, 8, 6\}$ (At least 4 levels).
    \end{itemize}
    
    \item \textbf{Quantization Strategy}
    \begin{itemize}
        \item PTQ vs. QAT (Main conclusions must rely on QAT) [43, 62].
        \item Uniform vs. Mixed-precision (Bound with Hardware Cost Model) [38, 39].
    \end{itemize}
    
    \item \textbf{Time Steps and Low-Latency Constraints}
    \begin{itemize}
        \item $T \in \{3, 4, 5\}$: Maintain consistency with EAS-SNN's low-step advantage and observe sensitivity of compression to low $T$ [1].
    \end{itemize}
\end{enumerate}

\subsubsection{A4) Output Forms (Design for Manuscript)}
\begin{itemize}
    \item \textbf{Table-A}: Gen1 mAP (SNN-FP / SNN-P / SNN-Q / SNN-Q(W+U) / SNN-HAPQ + ANN baselines).
    \item \textbf{Fig-A1}: Accuracy-Compression Rate Curve (mAP vs. $\kappa$ / bitwidth).
    \item \textbf{Fig-A2}: Membrane Potential Quantization Sensitivity (mAP vs. $b_U$; with fixed W bitwidth).
\end{itemize}

\subsection{B. FPGA: Comparative Experiments (Goal: Significant Performance Boost)}

\subsubsection{B1) FPGA Experimental Groups (Performance Bench)}
\textbf{Fixed Settings:} Identical FPGA board; Identical clock domain strategy; Identical DDR frequency; Identical input window (Same $T$, Same event pre-processing).

\begin{enumerate}
    \item \textbf{FPGA-Baseline}: SNN-FP (or ``FP16/INT16 Fixed-Point Implementation'' as hardware benchmark).
    \item \textbf{FPGA-P}: Structured Pruning Mapping only (Verify if sparsity yields real acceleration) [78].
    \item \textbf{FPGA-Q(W)}: Weight Quantization Mapping only (Verify DSP SIMD packing) [38, 39].
    \item \textbf{FPGA-Q(W+U)}: Weight + Membrane Quantization Mapping (Verify reduction in state bandwidth/BRAM pressure) [76, 77].
    \item \textbf{FPGA-HAPQ}: Final Joint Design (Main Conclusion).
\end{enumerate}

\begin{quote}
    \textit{Key Point: The FPGA comparison must include ``Q(W) vs. Q(W+U)''—otherwise, it is impossible to prove that membrane potential quantization is the solution to edge-side bottlenecks [76, 77].}
\end{quote}

\subsubsection{B2) FPGA Metrics (Directly Measured via Vivado/Board)}
\begin{itemize}
    \item \textbf{Resources}: DSP / BRAM / LUT / FF (Post-route).
    \item \textbf{Frequency}: $f_{clk}$ (Post-route timing).
    \item \textbf{Performance}: Latency (per window / per sample), Throughput (FPS or samples/s).
    \item \textbf{Power/Energy}: Board-level Power (W) + Energy per inference (J).
    \item \textbf{Bandwidth}: DDR Read/Write bandwidth usage (via AXI performance counters or ILA statistics).
\end{itemize}

\subsubsection{B3) Critical Comparisons (Essential Aspects for ``Significant Boost'')}
\begin{itemize}
    \item \textbf{Iso-Accuracy Comparison}: Compare Latency/Energy under mAP drop $\le$ threshold (e.g., 1\%) (Real benefit of compression).
    \item \textbf{Iso-Resource Comparison}: Compare maximum achievable mAP under the same DSP/BRAM limit (Highlighting ``Hardware-Awareness'') [38--41, 78].
    \item \textbf{Bottleneck Attribution}: Report Compute Cycles vs. State R/W Cycles (Prove that membrane quantization shifts the bottleneck structure) [76, 77].
\end{itemize}

\subsection{C. FPGA: Simulating Different Edge Compute Tiers (Goal: HAPQ Potential under Limited Compute)}

\subsubsection{C1) Definition of ``Compute Tiers'' (Hard Constraints)}
Define a budget vector for each tier:
\begin{equation}
    \mathcal{B}_k = \{ \tau^{(k)}_{DSP}, \tau^{(k)}_{BRAM}, \tau^{(k)}_{BW}, \tau^{(k)}_{fclk} \}
\end{equation}
\begin{itemize}
    \item \textbf{Tiny / Small / Medium / Large}: At least 4 tiers.
    \item $\tau_{BW}$ can be simulated by limiting AXI bursts or artificially lowering DDR frequency.
    \item $\tau_{fclk}$ can be simulated via clock constraints or inserting pipeline stages.
\end{itemize}

\subsubsection{C2) Method Comparison (Proving HAPQ Advantage)}
For each budget $\mathcal{B}_k$:
\begin{enumerate}
    \item \textbf{Naive Scaling}: Manual channel/width scaling (No hardware alignment).
    \item \textbf{Quant-only Policy}: Uniform bitwidth / HAQ-style (No state constraints) [38].
    \item \textbf{Prune-only Policy}: Magnitude/Channel pruning only (No hardware alignment) [48].
    \item \textbf{HAPQ}: NAS + Structured Pruning + W/U Quantization (With discrete resource modeling) [38--41, 78, 76, 77].
\end{enumerate}

\subsubsection{C3) Output Forms (Evidence Chain for ``Potential'')}
\begin{itemize}
    \item \textbf{Fig-C1 (Core)}: Pareto Front (mAP vs. Latency; mAP vs. Energy) colored by budget tiers.
    \item \textbf{Table-C}: The ``Best Feasible Design'' for each budget tier: $(mAP, Lat, Energy, DSP, BRAM, f_{clk})$.
    \item \textbf{Calib-Check}: Select at least 1--2 budget tiers for real on-board verification to calibrate the cost model (remaining tiers can use Post-route + Cycle model).
\end{itemize}

\section*{Minimum Executable Experiment Sequence (Immediate Actions)}
\textit{To avoid over-expanding initially:}
\begin{enumerate}
    \item \textbf{Gen1}: SNN-FP $\to$ SNN-Q(W) $\to$ SNN-Q(W+U) $\to$ SNN-HAPQ (Solidify the value of membrane potential quantization first) [1, 76, 77].
    \item \textbf{FPGA}: Implement only two versions initially (Baseline vs. HAPQ) to secure hard evidence on Resources/Frequency/Latency/Power [78].
    \item \textbf{Supplement Later}: Add Prune-only and Quant-only FPGA versions for attribution and ablation closure [38, 48].
\end{enumerate}

\end{document}